% Section 3.4, from lectures/L03/README.md: what you should be able to do, questions to test
% yourself, and the next lecture.

\section{Review}
\label{c3:sec:review}

\needspace{6\baselineskip}
\subsection*{What you should be able to do now}
\begin{itemize}
  \item Turn a vector name into the \code{.org} directive that puts a handler at it, in both word
    and byte addresses, and say which of the two \code{.org} wants and which the simulator
    reports.
  \item Say what the hardware saves on entry to a handler and what it does not, and write the
    prologue and epilogue that make up the difference.
  \item Explain why \code{SREG} has to be saved before anything that writes flags, and restored
    after.
  \item Configure a pin change interrupt, and say what it tells you and what it cannot.
  \item Identify a variable shared between a handler and the main loop, say how wide the window is
    in cycles, and close it.
  \item Predict how long your own handler blocks interrupts, and check the prediction.
  \item Say which parts of interrupt timing this book's simulator models faithfully and which it
    does not, and what you would use instead for the parts it does not.
\end{itemize}

\needspace{6\baselineskip}
\subsection*{Questions to test yourself}
\begin{enumerate}
  \item PCINT0 is vector 3. What goes after \code{.org}, and what would the same line have to say
    if this book used \code{avr-gcc}?
  \item An interrupt fires while a \code{ret} is executing. How many cycles pass before the
    handler's first instruction, at the earliest and at the latest?
  \item Your handler does \code{inc r24} and nothing else. Name two things wrong with it.
  \item Two buttons are on port B. One of them is pressed. Which vector fires, and how does the
    handler find out which button it was?
  \item A handler increments a 16-bit counter and the main loop reads it with two \code{lds}
    instructions. Give a value the main loop can read that the counter has never held, and say how
    many cycles wide the window is.
  \item Your handler measures 104 cycles. For how long are interrupts disabled, and why is the
    answer larger than 104?
\end{enumerate}

\needspace{6\baselineskip}
\subsection*{Looking ahead}
Chapter~4 is about data structures with no compiler to lay them out for you:
\begin{itemize}
  \item \code{X}, \code{Y} and \code{Z} in earnest: three pointers, and the displacement
    addressing that makes a struct field access one instruction.
  \item Where the seven bytes of an LED structure should actually live, and why
    \code{RAMEND + 1} is not it.
  \item The stack, now that an interrupt can arrive on top of two nested calls.
  \item Arrays of drivers, and a subroutine that walks one.
\end{itemize}
